\begin{sheetframe}{navy}{01｜線形代数}{連立方程式から固有値へ}

\subhead{行基本変形：拡大係数行列を上から追う}
\[
[A\mid b]=
\left[
\begin{array}{ccc|c}
1&1&1&6\\
2&-1&1&3\\
1&2&-1&2
\end{array}
\right]
\xrightarrow[\ R_3\leftarrow R_3-R_1\ ]
{\,R_2\leftarrow R_2-2R_1\,}
\left[
\begin{array}{ccc|c}
1&1&1&6\\
0&-3&-1&-9\\
0&1&-2&-4
\end{array}
\right]
\xrightarrow{\ \substack{
R_2\leftrightarrow R_3,\ R_3\leftarrow R_3+3R_2\\
R_3\leftarrow-\frac17R_3,\ R_2\leftarrow R_2+2R_3,\
R_1\leftarrow R_1-R_2-R_3}\ }
\left[
\begin{array}{ccc|c}
1&0&0&1\\
0&1&0&2\\
0&0&1&3
\end{array}
\right]
\quad\Longrightarrow\quad
(x,y,z)=(1,2,3).
\]
\[
\begin{array}{@{}ll@{\qquad}ll@{}}
[\,0\ \cdots\ 0\mid c\,],\ c\ne0
  &\Rightarrow \text{解なし},&
\text{自由変数が }r\text{ 個}
  &\Rightarrow \dim\ker A=r,\\
\text{先頭の1が }s\text{ 個}
  &\Rightarrow \rank A=s,&
\text{その位置に対応する元の列}
  &\Rightarrow \Img A\text{ の基底}.
\end{array}
\]
\[
\boxed{\dim\ker A+\rank A=n},
\qquad
\F=\R\text{ または }\C,\qquad
\ker A=\text{核},\quad \Img A=\text{像}.
\]

\vspace{.6mm}{\color{line}\rule{\textwidth}{.35pt}}\vspace{.5mm}

\begin{columns}[T,onlytextwidth]
\begin{column}{.318\textwidth}
\subhead{ベクトル空間：二つの演算と公理}
$\F=\R$ または $\C$。空でない集合 $V$ に二つの演算
\[
+:V\times V\to V,\qquad \cdot:\F\times V\to V,\quad(a,u)\mapsto au
\]
があり、$0\in V$ と各 $u\in V$ の $-u\in V$ が存在する。
任意の $u,v,w\in V$、$a,b\in\F$ に対し
\[
\begin{array}{@{}ll@{}}
u+v=v+u,&(u+v)+w=u+(v+w),\\
u+0=u,&u+(-u)=0,\\
a(u+v)=au+av,&(a+b)u=au+bu,\\
(ab)u=a(bu),&1u=u
\end{array}
\]
を満たすとき $V$ を $\F$ 上のベクトル空間という。

\minorhead{部分空間の定義（記号：$W\le V$）}
部分集合 $W\subset V$ が次を満たすこと：
\[
\boxed{W\ne\varnothing,\qquad
au+bv\in W
\quad(u,v\in W,\ a,b\in\F).}
\]

\minorhead{張る空間・一次独立・基底}
\[
\Span(v_1,\ldots,v_m)
=\left\{\sum_{i=1}^ma_iv_i:a_i\in\F\right\}.
\]
\[
v_1,\ldots,v_m\text{ が一次独立}
\iff
\sum_i a_iv_i=0\Rightarrow a_1=\cdots=a_m=0.
\]
一次独立かつ $\Span(v_1,\ldots,v_m)=V$ なら $V$ の基底。

\minorhead{線形写像 $f:V\to U$（$U$ も $\F$ 上のベクトル空間）}
\[
f(au+bv)=af(u)+bf(v)
\quad(u,v\in V,\ a,b\in\F).
\]
\[
\boxed{\dim V=\dim\ker f+\dim\Img f},
\qquad f\text{ は単射}\iff\ker f=\{0\}.
\]

\separator
\subhead{連立方程式と逆行列}
\[
Ax=b\text{ に解あり}\iff\rank A=\rank[A\mid b],
\]
\[
A^{-1}\exists\iff\det A\ne0\iff\rank A=n\iff\ker A=0.
\]

\separator
\subhead{特性多項式・対角化の判定}
\[
\chi_A(t)=\det(tI-A),\qquad \chi_A(A)=0
\quad\text{（Cayley--Hamilton）}.
\]
固有値を重複も含めて $\lambda_1,\ldots,\lambda_n$ とすると
\[
\tr A=\sum_{i=1}^n\lambda_i,\qquad
\det A=\prod_{i=1}^n\lambda_i.
\]
\[
\boxed{
A\text{ が対角化可能}
\iff
\sum_{\lambda}\dim\ker(A-\lambda I)=n.}
\]
\[
\dim\ker(A-\lambda I)
\le\lambda\text{ の特性多項式での重複度}.
\]
すべての固有値で等号なら対角化可能。特に相異なる固有値が $n$ 個あればよい。

\end{column}

\begin{column}{.356\textwidth}
\subhead{固有値：$A-\lambda I$ を行簡約する}
\[
A=
\begin{pmatrix}2&1\\1&2\end{pmatrix},
\qquad
\det(A-\lambda I)
=(2-\lambda)^2-1.
\]
\[
\det(A-\lambda I)=0
\iff(\lambda-1)(\lambda-3)=0
\iff\lambda=1,3.
\]
\[
\begin{aligned}
\lambda=1:\quad
A-I
&=\begin{pmatrix}1&1\\1&1\end{pmatrix}
\sim\begin{pmatrix}1&1\\0&0\end{pmatrix}\\
&\Rightarrow
E_1=\ker(A-I)
=\Span\!\left\{\binom{1}{-1}\right\},\\[1mm]
\lambda=3:\quad
A-3I
&=\begin{pmatrix}-1&1\\1&-1\end{pmatrix}
\sim\begin{pmatrix}1&-1\\0&0\end{pmatrix}\\
&\Rightarrow
E_3=\ker(A-3I)
=\Span\!\left\{\binom{1}{1}\right\}.
\end{aligned}
\]
この $A$ は実対称なので、異なる固有値の固有ベクトルは直交する。長さを $1$ にすると
\[
q_1=\frac1{\sqrt2}\binom{1}{-1},\qquad
q_2=\frac1{\sqrt2}\binom{1}{1},\qquad
Q=(q_1\ q_2)=\frac1{\sqrt2}
\begin{pmatrix}1&1\\-1&1\end{pmatrix}.
\]
\[
 Aq_1=q_1,\quad Aq_2=3q_2
\ \Longrightarrow\
AQ=Q\underbrace{\begin{pmatrix}1&0\\0&3\end{pmatrix}}_{D}
\ \Longrightarrow\
\boxed{Q^{\mathsf T}AQ=D},
\]
ここで $Q^{-1}=Q^{\mathsf T}$。これが正規直交対角化の流れ。

\separator
\subhead{Jordan 標準形：ブロックの形を読む}
\[
J_k(\lambda)=
\left(\begin{smallmatrix}
\lambda&1&&0\\
0&\lambda&\ddots&\\
\vdots&\ddots&\ddots&1\\
0&\cdots&0&\lambda
\end{smallmatrix}\right)_{k\times k}.
\]
対角が $\lambda$、その一段上が $1$、他が $0$。これを固有値 $\lambda$ の Jordan ブロックという。
\[
A\in M_n(\C):\quad
P^{-1}AP=\diag\bigl(J_{k_1}(\lambda_1),\ldots,J_{k_s}(\lambda_s)\bigr).
\]

$N=A-\lambda I$ とおくと
\[
\begin{aligned}
\text{ブロック（小行列）の個数}&=\dim\ker N,\\
\text{大きさが }j\text{ 以上の個数}
  &=\dim\ker N^j-\dim\ker N^{j-1}.
\end{aligned}
\]
\[
Nv=0,\qquad Nu=v
\quad\Longrightarrow\quad (v,u)\text{ の順で }J_2(\lambda)\text{ を得る}.
\]

\separator
\subhead{内積・スペクトル定理}
\[
A^{\mathsf T}=A\text{ または }A^*=A
\Longrightarrow
\begin{matrix}\text{固有値は実数},\\[-.3mm]\text{正規直交固有基底あり}.
\end{matrix}
\]
\end{column}

\begin{column}{.302\textwidth}
\subhead{射影：一つのベクトルを二方向へ分ける}
\begin{center}
\begin{tikzpicture}[x=1mm,y=1mm]
  \draw[muted,line width=.45pt] (-20,0)--(22,0) node[right,tinylabel] {$\Img P$};
  \draw[muted,line width=.45pt] (0,-4)--(0,20) node[above,tinylabel] {$\ker P$};
  \draw[flow] (0,0)--(16,13) node[midway,above left,tinylabel] {$x$};
  \draw[flow] (0,0)--(16,0) node[midway,below,tinylabel] {$Px$};
  \draw[flow] (16,0)--(16,13) node[midway,right,tinylabel] {$(I-P)x$};
\end{tikzpicture}
\end{center}
\[
P^2=P\Rightarrow
V=\Img P\oplus\ker P,\qquad x=Px+(I-P)x.
\]
ベクトル $u\ne0$ の張る直線への直交射影は
\[
\boxed{P=\frac{uu^*}{u^*u},\qquad Px=\frac{u^*x}{u^*u}\,u.}
\]

\separator
\subhead{不変平面：法線を見る}
\begin{center}
\begin{tikzpicture}[x=1mm,y=1mm]
  \fill[accent!7] (-20,-3)--(13,-3)--(20,7)--(-13,7)--cycle;
  \draw[accent,line width=.5pt] (-20,-3)--(13,-3)--(20,7)--(-13,7)--cycle;
\draw[flow] (0,2)--(0,17) node[above,tinylabel] {$u$};
  \node[tinylabel] at (-13,1) {$W=\ker u^{\mathsf T}$};
  \node[tinylabel] at (14,12) {$A^{\mathsf T}u=\lambda u$};
\end{tikzpicture}
\end{center}
実対称行列では
\[
\lambda_{\min}\le\frac{x^{\mathsf T}Ax}{x^{\mathsf T}x}\le\lambda_{\max}
\qquad(x\ne0).
\]
$A=A^{\mathsf T}$、$A_k=$ 左上 $k\times k$ 部分行列なら
\[
x^{\mathsf T}Ax>0\ (x\ne0)
\iff \text{全固有値が正}
\iff \det A_k>0\ (k=1,\ldots,n).
\]

\separator
\subhead{Gram--Schmidt の直交化}
\[
u_k=v_k-\sum_{j<k}\frac{\ip{u_j}{v_k}}{\ip{u_j}{u_j}}u_j,
\qquad e_k=\frac{u_k}{\norm{u_k}}.
\]
\end{column}
\end{columns}
\end{sheetframe}
